% Secciones 1--3 del delta gravitatorio íntegro.
\section{Domaine généalogique de la réalisation gravitationnelle}
\label{sec:l9s102:gravedad-dominio-genealogico}

La réalisation gravitationnelle reçoit comme objet de départ une section du
continu discret déjà construite par la composition
\begin{equation}
 \mathrm{APP}\longrightarrow\mathrm{TRIT}\longrightarrow\mathrm{TPK}
 \longrightarrow X_{\mathrm{TPK}}^{\mathrm{enr}}
 \longrightarrow\mathcal C_{\mathrm{cont}}^{\mathrm{disc}}.
 \label{eq:l9s102:gravedad-cadena-genealogica}
\end{equation}
APP conserve les évaluations additive et multiplicative au moyen du résidu,
du quotient et de la retenue ; TRIT détermine le régime et l'orientation ; le TPK
incorpore le sélecteur tritique de phase sans éliminer aucun des deux feuillets
APP, effectue le transport cardinal, conserve la mémoire et réalise le retour
de phase sans réinitialiser l'état. Le transducteur
dimensionnel agit ensuite sur la section survivante du continu :
\begin{equation}
 \mathfrak G_{108}:\mathcal C_{\mathrm{cont}}^{\mathrm{disc}}
 \times\mathcal L_S^+
 \longrightarrow\mathcal G_{\mathrm{HMT\text{-}MD}},
 \qquad
 (x,\hbar_a)\longmapsto
 \bigl(G_a,R_{108},m_{108,a},\ell_{P,a}^{,2}\bigr).
 \label{eq:l9s102:gravedad-mapa-realizacion}
\end{equation}
Les cardinaux \(90/120\), la périodicité \(108\), l'action et le couple
angulaire entrent avec leur généalogie conservée. Aucune identification
gravitationnelle ultérieure ne sélectionne rétroactivement ces antécédents.

\section{Transducteur gravitationnel et sélecteur chronologique de périodicité}
\label{sec:l9s102:gravedad-transductor}

Soient \(c_{\mathrm{int}}>0\), \(t_0>0\),
\(\ell_0=c_{\mathrm{int}}t_0\) et
\(\hbar_a\in\mathcal L_S^+\), avec
\(a\in\{\mathrm{pre},\mathrm{ret}\}\). Le transducteur dimensionnel est
\begin{equation}
 G_a(\theta)
 =\theta\frac{c_{\mathrm{int}}^5t_0^2}{\hbar_a}
 =\theta\frac{c_{\mathrm{int}}^3\ell_0^2}{\hbar_a},
 \qquad \theta>0.
 \label{eq:l9s102:g-transductor}
\end{equation}
La chronologie cyclique de \(108\) états détermine
\begin{equation}
 T_{108}=108t_0,
 \qquad
 \omega_{108}=\frac{2\pi_{\rm HMT}}{108t_0},
 \qquad
 R_{108}=\frac{c_{\mathrm{int}}}{\omega_{108}}
 =\frac{54}{\pi_{\rm HMT}}\ell_0.
 \label{eq:l9s102:g-reloj108}
\end{equation}
Le quantum chronologique et sa masse associée sont
\begin{equation}
 E_{108,a}=\hbar_a\omega_{108},
 \qquad
 m_{108,a}=\frac{E_{108,a}}{c_{\mathrm{int}}^2}.
 \label{eq:l9s102:g-cuanto-masa}
\end{equation}
La longueur de Compton réduite satisfait
\begin{equation}
 \bar\lambda_{C,a}
 =\frac{\hbar_a}{m_{108,a}c_{\mathrm{int}}}=R_{108}.
 \label{eq:l9s102:g-compton}
\end{equation}

\begin{theorem}[Unicité du sélecteur gravitationnel de périodicité]
\label{thm:l9s102:g-selector}
Avec la convention réduite \(r_g=Gm/c_{\mathrm{int}}^2\), la condition
\(r_{g,a}(\theta)=R_{108}\) possède l'unique solution positive
\begin{equation}
 \boxed{
 \theta_{108}^{\mathrm{circ}}=\left(\frac{54}{\pi_{\rm HMT}}\right)^2.}
 \label{eq:l9s102:g-theta}
\end{equation}
La réalisation correspondante identifie structurellement
\begin{equation}
 R_{108}=\bar\lambda_{C,a}=r_{g,a}=\ell_P^{\mathrm{circ}}.
 \label{eq:l9s102:g-cuatro-longitudes}
\end{equation}
\end{theorem}

\begin{proof}
Par \cref{eq:l9s102:g-transductor,eq:l9s102:g-cuanto-masa},
\[
 r_{g,a}(\theta)
 =\frac{G_a(\theta)m_{108,a}}{c_{\mathrm{int}}^2}
 =\theta\frac{\pi_{\rm HMT}}{54}\ell_0.
\]
L'égalité avec \(R_{108}=(54/\pi_{\rm HMT})\ell_0\) impose
\(\theta=(54/\pi_{\rm HMT})^2\). La linéarité en \(\theta\) et la positivité des
facteurs déterminent l'unicité.
\end{proof}

Dans la carte où \(c_{\mathrm{int}}=(\mu\varepsilon)^{-1/2}\), le
transducteur adopte la forme
\begin{equation}
 \boxed{
 G_a=\left(\frac{54}{\pi_{\rm HMT}}\right)^2
 \frac{t_0^2}{\hbar_a(\mu\varepsilon)^{5/2}}.}
 \label{eq:l9s102:g-mu-epsilon}
\end{equation}
Le produit constitutif \(\mu\varepsilon\) gouverne le cône causal et
l'échelle gravitationnelle ; le quotient \(Z^2=\mu/\varepsilon\) conserve
l'orientation d'impédance.

\section{Covariance bidirectionnelle de la gravité et de l'action}
\label{sec:l9s102:g-hbar}

On définit
\begin{equation}
 R_{\mathrm{act}}=\frac{\hbar_{\mathrm{pre}}}
 {\hbar_{\mathrm{ret}}}.
 \label{eq:l9s102:r-act}
\end{equation}
La masse chronologique et le couplage gravitationnel se transforment
de manière contragrédiente :
\begin{equation}
 \frac{m_{108,\mathrm{pre}}}{m_{108,\mathrm{ret}}}=R_{\mathrm{act}},
 \qquad
 \frac{G_{\mathrm{pre}}}{G_{\mathrm{ret}}}=R_{\mathrm{act}}^{-1}.
 \label{eq:l9s102:g-hbar-reciprocidad}
\end{equation}
Leur produit demeure fixe :
\begin{equation}
 \boxed{
 G_a\hbar_a
 =\theta_{108}^{\mathrm{circ}}c_{\mathrm{int}}^5t_0^2,
 \qquad
 \ell_P^2=\frac{G_a\hbar_a}{c_{\mathrm{int}}^3}
 =\theta_{108}^{\mathrm{circ}}\ell_0^2.}
 \label{eq:l9s102:g-hbar-invariante}
\end{equation}
L'action conserve l'orientation du feuillet ; l'aire de Planck conserve la
géométrie commune aux deux orientations.

\begin{theorem}[Indépendance de carte des sections \(h\) et \(G\)]
\label{thm:l9s102:covariancia-unidades-g}
Soient \(u_L,u_T,u_S\) des bases positives des droites dimensionnelles de
longueur, de temps et d'action, et remplaçons-les par
\[
 u_L'=\lambda_Lu_L,\qquad
 u_T'=\lambda_Tu_T,\qquad
 u_S'=\lambda_Su_S,
 \qquad \lambda_L,\lambda_T,\lambda_S>0.
\]
Les coordonnées des mêmes sections vérifient
\begin{equation}
 c_{\rm int}'=\frac{\lambda_T}{\lambda_L}c_{\rm int},\qquad
 t_0'=\frac{t_0}{\lambda_T},\qquad
 \hbar_a'=\frac{\hbar_a}{\lambda_S},\qquad
 G_a'=\frac{\lambda_S\lambda_T^3}{\lambda_L^5}G_a.
 \label{eq:l9s102:cambio-carta-g}
\end{equation}
Pour la base induite
\(u_G=u_L^5u_T^{-3}u_S^{-1}\), on obtient
\begin{equation}
 \hbar_a'u_S'=\hbar_au_S,
 \qquad G_a'u_G'=G_au_G.
 \label{eq:l9s102:secciones-h-g-invariantes}
\end{equation}
Par conséquent,
\(h_a=2\pi_{\rm HMT}\hbar_a\),
\(G_a=\theta_{108}^{\rm circ}c_{\rm int}^5t_0^2/\hbar_a\) et
\(G_a\hbar_a=\theta_{108}^{\rm circ}c_{\rm int}^5t_0^2\)
sont des identités de sections, non des égalités dépendant d'un choix
d'unités. La carte métrologique publie leurs coordonnées numériques après
la construction HMT--MD.
\end{theorem}

\begin{proof}
Les trois premières transformations de
\eqref{eq:l9s102:cambio-carta-g} sont la loi contragrédiente des coordonnées
dans \(L\otimes T^{-1}\), \(T\) et \(S\). Leur substitution dans
\eqref{eq:l9s102:g-transductor} produit la quatrième. La base de la droite
gravitationnelle se transforme par le facteur inverse, ce qui prouve
\eqref{eq:l9s102:secciones-h-g-invariantes}. Les scalaires
\(\pi_{\rm HMT}\) et \(\theta_{108}^{\rm circ}\) sont adimensionnels et
demeurent invariants.
\end{proof}
